\section{Retracts and a constructed nondeterministic domain}
\label{sec:p10-domain-constructions-retractions}

The chain bilimit and canonical recursion arguments of
Theorems~\ref{thm:p10-domain-bilimit} and
\ref{thm:p10-domain-invariance-mixed} already construct recursive objects.
Here the additional tasks are to split continuous idempotents, to distinguish
retracts from subdomains, and to construct a free domain for finite
nondeterministic choice. The source is the corrected and expanded
Abramsky--Jung author presentation, Sections~3.1 and~6.1.2--6.2.1
\cite[pp.~34--41,87--95]{AbramskyJungCorrectedExpandedDomainTheory}.
Its printed chapter and the Gunter--Scott, Gunter and Mosses presentations
use their own conventions; equivalence with them requires a matched comparison.

\subsection{Retractions are mathematical representation data}

\begin{hbdefinition}{def:p10-dc-retraction}{Sections, retractions and idempotent images}
A continuous section--retraction pair consists of Scott-continuous
\(s:D\to E\), \(r:E\to D\) with \(rs=1_D\).
It need not satisfy \(sr\sqsubseteq1_E\); that additional inequality
makes it the embedding--projection pair of
Definition~\ref{def:p10-domain-expanding-sequence}.
A continuous idempotent \(f:E\to E\) satisfies \(f^2=f\).
Give \(\operatorname{im}(f)=\{x:f(x)=x\}\) its inherited order.
All objects and bases here are set-sized under the existing universe
convention; directed sets are nonempty. A section--retraction pair is
not by itself an effective encoding of a programming language
\cite[Definitions~3.1.1,3.1.15, pp.~34,39]
{AbramskyJungCorrectedExpandedDomainTheory}.
\end{hbdefinition}

\begin{hbtheorem}{thm:p10-dc-splitting}{proved=annex:p10.dc.splitting}{Actual splitting, topology and directed idempotents}
The image of a continuous idempotent on a dcpo is a dcpo with
directed suprema computed in its parent. Its inclusion \(s\) and
corestriction \(r(x)=f(x)\) are continuous and satisfy \(rs=1\),
\(sr=f\). Its Scott topology is the subspace topology inherited from
the parent. A directed family of continuous idempotents, ordered
pointwise, has a continuous idempotent supremum.
If \(E\) is continuous with basis \(B\), every continuous retract
\(D\) is continuous with basis \(r(B)\); countability of a basis is
preserved. This does not say that \(r\) preserves every way-below pair
or that algebraicity is preserved
\cite[Lemma~3.1.3; Theorem~3.1.4; Propositions~3.1.7--8, pp.~34--36]
{AbramskyJungCorrectedExpandedDomainTheory}.
\end{hbtheorem}

\begin{hbtheorem}{thm:p10-dc-algebraic-cover}{proved=annex:p10.dc.algebraic-cover}{Every continuous domain has an algebraic retract presentation}
Let \(B\) be a basis of a continuous domain \(D\). In this theorem
\(\operatorname{Idl}(B,\sqsubseteq)\) uses the restricted information
order, not the usually nonreflexive way-below relation. Put
\[
 s(x)=B_x=\{b\in B:b\ll x\},\qquad
 r(J)=\bigsqcup^{\uparrow}J.
\]
These are continuous maps and \(rs=1_D\).
The ideal domain is algebraic, with compact principal ideals
\(\downarrow_B b\). A countable \(B\) gives an omega-algebraic
cover of an omega-continuous \(D\).
For \(D=[0,1]\) and its rational basis, this is an algebraic domain
retracting onto a nonalgebraic continuous domain. Thus ordinary continuous
retractions genuinely fail to preserve algebraicity
\cite[Theorem~3.1.5; Proposition~3.1.6, p.~35]
{AbramskyJungCorrectedExpandedDomainTheory}.
\end{hbtheorem}

\begin{hbtheorem}{thm:p10-dc-function-retract}{proved=annex:p10.dc.function-retract}{Retraction transport for continuous functions}
For section--retraction pairs \((s_D,r_D)\), \((s_E,r_E)\),
from \(D,E\) into \(D',E'\), respectively, the maps
\[
 S(h)=s_Ehr_D,
 \quad R(k)=r_Eks_D,
 \quad [D\longrightarrow E]\mathrel{\mathop{\rightleftarrows}^{S}_{R}}
       [D'\longrightarrow E']
\]
are continuous and satisfy \(RS=1\).
If the original pairs are embedding--projection pairs, so is \((S,R)\).
Products transport retractions coordinatewise as well.
In particular this gives an actual retract representation of a function
domain, not merely a bijection of function names. Computability of those
names and maps still requires effective basis and representation data.
This is an explicit construction using the source's retraction and
function-space machinery \cite[pp.~34--36,43]
{AbramskyJungCorrectedExpandedDomainTheory}.
\end{hbtheorem}

\subsection{The Plotkin powerdomain from its approximation basis}

\begin{hbdefinition}{def:p10-dc-plotkin-basis}{Two-sided approximation and formal union}
Let \(D\) be a nonempty continuous domain with basis \(B\), and let
\(\mathcal P_f(B)\) mean its \emph{nonempty finite} subsets.
The source's Egli--Milner approximation relation is
\[
 M\ll_{\mathrm{EM}}N\quad\Longleftrightarrow\quad
 (\forall a\in M)(\exists b\in N)\ a\ll b
 \ \text{and}\
 (\forall b\in N)(\exists a\in M)\ a\ll b.
\]
Define \(P^{\mathrm P}(D)=
\operatorname{Idl}(\mathcal P_f(B),\ll_{\mathrm{EM}})\), using round
ideals as in Definition~\ref{def:p10-domain-abstract-basis}.
Write \(i(M)=\{L:L\ll_{\mathrm{EM}}M\}\) and
\[
 \eta_D(x)=\{M\in\mathcal P_f(B):(\forall b\in M)\ b\ll x\}.
\]
The symbol \(\cup\) has two expressly typed readings below.
For finite basis sets \(M,N\), it is ordinary set union.
For powerdomain ideals or elements of a dcpo-semilattice, it is
\emph{formal union}: a continuous commutative, associative, idempotent
binary operation. It is not assumed to be a join in the information order.
On powerdomain ideals define the latter operation by
\[
 I\cup J=\{L:(\exists M\in I)(\exists N\in J)\
                         L\ll_{\mathrm{EM}}(M\cup N)\}.
\]
It generally differs from their ordinary union as sets of finite sets
\cite[Definitions~6.2.1--2; Theorem~6.2.3, pp.~93--94]
{AbramskyJungCorrectedExpandedDomainTheory}.
\end{hbdefinition}

\begin{hbtheorem}{thm:p10-dc-plotkin-object}{proved=annex:p10.dc.plotkin-object}{Existence, approximation, and the generator map}
\((\mathcal P_f(B),\ll_{\mathrm{EM}})\) is an abstract basis.
Consequently \(P^{\mathrm P}(D)\) is a continuous dcpo with directed
unions as suprema and basis \(i(\mathcal P_f(B))\).
The map \(\eta_D:D\to P^{\mathrm P}(D)\) is Scott-continuous.
If \(D\) is algebraic, choosing \(B=K(D)\) makes the powerdomain
algebraic. Its compact elements are the principal ideals of the
Egli--Milner preorder, with equal principal ideals identified.
Countable bases yield countable powerdomain bases. If \(D\) is pointed
and \(\bot_D\in B\), then \(i(\{\bot_D\})\) is its bottom and
equals \(\eta_D(\bot_D)\)
\cite[Theorem~6.2.3; pp.~94--95]
{AbramskyJungCorrectedExpandedDomainTheory}.
\end{hbtheorem}

\begin{hbtheorem}{thm:p10-dc-formal-union}{proved=annex:p10.dc.formal-union}{A genuine continuous semilattice operation}
The displayed formal union produces round ideals and is jointly
Scott-continuous. It obeys commutativity, associativity and idempotence.
For every \(M\in\mathcal P_f(B)\),
\[
 i(M)=\mathop{\bigcup}_{b\in M}^{\mathrm{formal}}\eta_D(b).
\]
Here the finite formal union is unambiguous by the stated equations;
no nullary union or empty nondeterministic choice has been introduced.
Thus \(P^{\mathrm P}(D)\) is a dcpo-semilattice in the source's
sense \cite[pp.~93--94]{AbramskyJungCorrectedExpandedDomainTheory}.
\end{hbtheorem}

\begin{hbtheorem}{thm:p10-dc-free-powerdomain}{proved=annex:p10.dc.free-powerdomain}{The complete free universal property and its continuity}
Let \(A\) be any dcpo with a jointly Scott-continuous commutative,
associative, idempotent operation \(\cup_A\), and let \(g:D\to A\)
be Scott-continuous. There is exactly one continuous semilattice
homomorphism \(\operatorname{ext}(g):P^{\mathrm P}(D)\to A\)
with \(\operatorname{ext}(g)\eta_D=g\). It is explicitly
\[
 \operatorname{ext}(g)(I)=
 \bigsqcup^{\uparrow}_{M\in I}
                 \mathop{\bigcup}_{b\in M}^{A}g(b).
\]
The family inside the supremum is directed in \(A\)'s information
order. The assignment \(g\mapsto\operatorname{ext}(g)\) preserves
directed suprema in the pointwise function orders.
The target \(A\) need not itself be a continuous domain.
This is the Plotkin specialization of the source's free
continuous-domain algebra theorem, not a claim that every semantic
choice structure is a set of possible terminating outcomes
\cite[Theorems~6.1.6--7; Proposition~6.1.8; Theorem~6.2.3,
pp.~89--91,94]{AbramskyJungCorrectedExpandedDomainTheory}.
\end{hbtheorem}

\begin{hbtheorem}{thm:p10-dc-powerdomain-functor}{proved=annex:p10.dc.powerdomain-functor}{Functorial action, retract preservation and flattening}
For continuous \(u:D\to E\) between nonempty continuous domains put
\(P^{\mathrm P}(u)=\operatorname{ext}(\eta_Eu)\).
This is a locally continuous functor. It transports continuous
section--retraction pairs and embedding--projection pairs without
changing their inequalities. The maps
\[
 \mu_D=\operatorname{ext}(1_{P^{\mathrm P}(D)}):
       P^{\mathrm P}(P^{\mathrm P}(D))\to P^{\mathrm P}(D)
\]
and \(\eta_D\) are natural and satisfy
\[
 \mu_D\eta_{P^{\mathrm P}(D)}=1,
 \quad\mu_DP^{\mathrm P}(\eta_D)=1,
 \quad\mu_DP^{\mathrm P}(\mu_D)=\mu_D\mu_{P^{\mathrm P}(D)}.
\]
Thus the generator and flattening form a monad with an actual free
universal property. This conclusion does not assert closure of every
cartesian closed subclass of domains under the powerdomain constructor
\cite[Theorem~6.1.9; pp.~91,93--95]
{AbramskyJungCorrectedExpandedDomainTheory}.
\end{hbtheorem}

\begin{hbtheorem}{thm:p10-dc-choice-example}{proved=annex:p10.dc.choice-example}{A finite nondeterministic example and failed substitutes}
Take the flat finite domain \(D=\{\bot,a,b\}\), with incomparable
\(a,b\). The seven nonempty subsets represent seven distinct compact
powerdomain elements, ordered by \(\sqsubseteq_{\mathrm{EM}}\).
In particular
\[
 \{\bot\}\sqsubseteq_{\mathrm{EM}}\{\bot,a\}
 \sqsubseteq_{\mathrm{EM}}\{a\},\qquad
 \{a\}\not\sqsubseteq_{\mathrm{EM}}\{a,b\},
 \quad\{a,b\}\not\sqsubseteq_{\mathrm{EM}}\{a\}.
\]
Formal union of the represented outcomes \(\{a\}\), \(\{b\}\)
is \(\{a,b\}\), but that element is not an information-order
upper bound of either input. Ordinary union of their principal ideals
is not directed and is not a powerdomain ideal.
If empty finite subsets are admitted without changing the theory,
the empty set has no Egli--Milner relation to a nonempty set and supplies
an extra isolated alternative, not the existing bottom \(\{\bot\}\).
Finally an ordinary retract can create compactness: the constant map
to \(1/2\) on \([0,1]\) has a compact one-point image although
\(1/2\not\ll1/2\) in the parent.
These are explicit construction/failure cases, not operational
adequacy claims for a programming language
\cite[pp.~35--36,93--96]{AbramskyJungCorrectedExpandedDomainTheory}.
\end{hbtheorem}

\subsection{An explicit effective algebraic regime}

\begin{hbdefinition}{def:p10-dc-effective}{Numbered compact bases and effective maps}
In this subsection an object is a nonempty algebraic domain with a
fixed numbering \(K(D)=\{b_n:n\in\mathbb N\}\), allowing repeated
names, whose compact order \(b_i\sqsubseteq b_j\) is recursive in
the indices. The numbering and decision procedure are part of the object.
For a point \(x\), put \(I_x=\{n:b_n\sqsubseteq x\}\); call
\(x\) computable when \(I_x\) is recursively enumerable.
For a continuous \(f:D\to E\), with target compact names \(c_m\),
put
\[
 R_f=\{(n,m):c_m\sqsubseteq f(b_n)\};
\]
call \(f\) effective when this relation is recursively enumerable.
This is an explicit algebraic specialization of the source's numbered
abstract bases and recursively enumerable approximable mappings
\cite[Section~8.1.1, p.~145]{AbramskyJungCorrectedExpandedDomainTheory}.
It does not identify all Scott-continuous maps with computable maps.
All finite choice codes below are nonempty finite lists of compact
indices; repetition merely gives alternative names for the same finite
basis set. No algorithm deciding arbitrary program termination is assumed.
\end{hbdefinition}

\begin{hbtheorem}{thm:p10-dc-effective-choice}{proved=annex:p10.dc.effective-choice}{An actual effective category and effective choice construction}
These objects and effective continuous maps form a category. Effective
maps send computable points to computable points. The Plotkin powerdomain
has an effective compact-basis numbering by finite choice codes:
its compact order is the recursive finite Egli--Milner comparison.
The functor \(P^{\mathrm P}\) preserves effective maps, and its unit
and flattening maps are effective in these numberings. In particular the
monad of Theorem~\ref{thm:p10-dc-powerdomain-functor} restricts to this
explicit effective category.
Continuity alone is insufficient: if \(U\subseteq\mathbb N\) is
not recursively enumerable, the constant map from the one-point domain
to \(\mathcal P(\mathbb N)\) with value \(U\) is continuous but not
effective in the standard finite-subset compact numbering.
These are proved finite-basis lifts of the constructed choice model,
not a claim to have supplied every effective function-space, fixed-point
or universal-domain result mentioned by the source's overview
\cite[pp.~94,145]{AbramskyJungCorrectedExpandedDomainTheory}.
\end{hbtheorem}

\begin{hbmethod}{def:p10-dc-using-the-model}{Using retractions and finite-choice semantics}
Specify the source and target dcpo's, their bases, and the actual maps
\(s,r\); prove \(rs=1\) and, when a subdomain is claimed,
\(sr\sqsubseteq1\). For finite nondeterministic observations, specify
the two-sided approximation relation before forming round ideals.
Interpret a generator by a continuous \(g\), evaluate finite formal
choices with \(\cup_A\), and use the proved directed extension for
infinite information. Keep information refinement distinct from
semilattice accumulation. An effective model additionally needs
enumerated basis codes and computable approximation/interpretation maps;
an operational adequacy theorem additionally needs a specified language,
reduction, observation and logical relation. None of these data follows
from the free universal property alone.
\end{hbmethod}

\begin{proofboundary}{def:p10-dc-source-boundary}{What this construction does and does not establish}
The preceding proofs supply retraction representation and a free
continuous nondeterministic-domain construction. They do not establish
an operational/denotational adequacy or full-abstraction theorem,
the general effective continuous-domain/function-space/recursive-equation
theory beyond the declared algebraic choice regime, the inverse-limit theorem for directed-index
adjunction diagrams, or the full bifinite/FS-domain hierarchy.
The further convex-hull and topological representations of the Plotkin
construction, one-sided powerdomains and probabilistic powerdomains
require their own hypotheses and proofs.
The distinct Gunter--Scott, Mosses and type-semantics presentations are
not identified with this corrected author presentation. The results here
apply to the specified continuous-domain and effective algebraic regimes,
not to every domain model or programming-language interpretation.
\end{proofboundary}
